% Rewritten theory for Lab 07: BJT biasing (fixed and self-bias)
\section{Theory}
\markright{Theory}

This experiment studies the DC operating (quiescent) point of a bipolar junction transistor (BJT) and the stability
of that operating point with respect to variations in transistor current gain $\beta$. A stable Q-point keeps the BJT in
active mode despite changes in temperature or device-to-device $\beta$ variation, which reduces amplifier distortion
and improves predictable gain \cite{boylestad2007,sedra1991}.

\subsection{Bias network and Thevenin equivalent}
The common bias used here is the single-supply voltage divider feeding the base, with collector and emitter resistors
shown in Figure~\ref{lab07_fig:bias_circuit} (placeholder). The divider (resistors $R_{B1}$ and $R_{B2}$) can be
replaced by its Thevenin equivalent, with
\begin{equation}
  V_{TH}=\frac{R_{B2}}{R_{B1}+R_{B2}}\,V_{CC},\qquad R_{TH}=\frac{R_{B1}R_{B2}}{R_{B1}+R_{B2}}.
\end{equation}

Applying KVL at the base loop and using $I_E\approx(1+\beta)I_B$ gives the base current
\begin{equation}
  I_B=\dfrac{V_{TH}-V_{BE}}{R_{TH}+(1+\beta)R_E},
\end{equation}
where $V_{BE}\simeq0.7\,$V in the active region. The quiescent collector and emitter currents follow:
\begin{align}
  I_{C_Q} & =\beta\,I_B,     \\
  I_{E_Q} & =(1+\beta)\,I_B,
\end{align}
and the collector–emitter voltage at Q is
\begin{equation}
  V_{CE_Q}=V_{CC}-I_{C_Q}R_C-I_{E_Q}R_E.
\end{equation}

These closed-form expressions allow direct evaluation of how $I_C$, $I_B$, and $V_{CE}$ change when $\beta$
varies, so the stability of a given bias network can be compared analytically \cite{keown2001}.

\begin{figure}[H]
  \centering
  \fbox{\parbox[c][3cm][c]{0.85\linewidth}{\centering Placeholder for Figure: bias network and Thevenin equivalent}}
  \caption{Bias network and its Thevenin equivalent.}
  \label{lab07_fig:bias_circuit}
\end{figure}

\subsection{Fixed bias vs self-bias}
The two practical circuits used in this laboratory are the fixed-bias circuit and the self-bias (emitter-feedback) circuit
shown schematically in Figure~\ref{lab07_fig:fixed_self} (placeholder). In the fixed-bias arrangement the base is driven
through a single resistor (or potentiometer), so the base current depends strongly on $\beta$ and temperature. In the
self-bias network the emitter resistor provides negative feedback: when $I_C$ increases, $V_E$ rises, which reduces
base–emitter forward bias and tends to restore the operating point. This feedback makes self-bias circuits much more
stable against $\beta$ variation and parameter drift \cite{comer2002}.

For amplifier design it is common to choose component values so that
\begin{equation}
  V_{CE_Q}\approx\frac{V_{CC}}{2}
\end{equation}
to maximise symmetrical output swing without clipping.

\begin{figure}[H]
  \centering
  \fbox{\parbox[c][3cm][c]{0.85\linewidth}{\centering Placeholder for Figure: fixed-bias and self-bias circuit diagrams}}
  \caption{Schematic placeholders for (a) fixed-bias and (b) self-bias CE circuits.}
  \label{lab07_fig:fixed_self}
\end{figure}
